\documentclass[11pt,a4paper]{article}

\usepackage[T1]{fontenc}
\usepackage[utf8]{inputenc}
\usepackage{lmodern}
\usepackage{microtype}
\usepackage[a4paper,margin=2.5cm]{geometry}
\usepackage{amsmath,amssymb,mathtools}
\usepackage{bm}
\usepackage{booktabs}
\usepackage{array}
\usepackage{xcolor}
\usepackage{hyperref}

\hypersetup{
  colorlinks=true,
  linkcolor=blue!50!black,
  citecolor=blue!50!black,
  urlcolor=blue!50!black
}

\newcommand{\E}{\mathbf E}
\newcommand{\Hfield}{\mathbf H}
\newcommand{\Bfield}{\mathbf B}
\newcommand{\Dfield}{\mathbf D}
\newcommand{\J}{\mathbf J}
\newcommand{\Pfield}{\mathbf P}
\newcommand{\nvec}{\hat{\mathbf n}}
\newcommand{\dd}{\mathrm d}
\newcommand{\ee}{\mathrm e}
\newcommand{\ii}{\mathrm i}
\newcommand{\eps}{\varepsilon}
\newcommand{\epsinf}{\varepsilon_\infty}
\newcommand{\epsd}{\varepsilon_d}
\newcommand{\epst}{\varepsilon_T}
\newcommand{\epsl}{\varepsilon_L}
\newcommand{\omegap}{\omega_p}
\newcommand{\vF}{v_F}

\title{\textbf{The vanilla hard-wall hydrodynamic Drude model}\\
\large A derivation with physical interpretation}
\author{}
\date{27 July 2026}

\begin{document}

\maketitle

\begin{abstract}
This note derives the standard linearized hard-wall hydrodynamic Drude model
(HDM) for a metal and explains what each equation means physically.  The
central modification of local Drude theory is the finite compressibility of the
conduction-electron fluid.  It produces a pressure term proportional to
\(\nabla(\nabla\!\cdot\!\J)\), leaves the transverse response local, and gives
the longitudinal electron-density response a finite wave number.  The note
also distinguishes the incident electromagnetic field from the total
self-consistent field, explains why local Drude theory already supports surface
plasmons, derives the electric-field-only formulation by eliminating \(\J\),
and derives the additional hard-wall boundary condition \(\nvec\!\cdot\!\J=0\).
\end{abstract}

\tableofcontents

\section{Scope, notation, and the central idea}

We consider a metal occupying a region \(\Omega_m\), surrounded by a local
dielectric with relative permittivity \(\epsd\).  The unit normal \(\nvec\)
points from the metal into the dielectric.  The metal contains a mobile
conduction-electron fluid.  Bound electrons and interband processes are
represented by a local background permittivity \(\epsinf\).  The mobile
electrons are treated hydrodynamically.

Throughout the derivation we use the harmonic convention
\begin{equation}
    \text{all fields}(\mathbf r,t)=\Re\!\left[\text{field}(\mathbf r,\omega)
    \ee^{-\ii\omega t}\right].
    \label{eq:time-convention}
\end{equation}
With this convention,
\begin{equation}
    \partial_t\longrightarrow-\ii\omega.
    \label{eq:dt-rule}
\end{equation}

The most important conceptual point is the following:
\begin{equation}
    \boxed{\text{The HDM equations inside the metal act on the total
    self-consistent field }\E.}
    \label{eq:total-field-box}
\end{equation}
For a scattering problem one may write
\begin{equation}
    \E_{\mathrm{tot}}=\E_{\mathrm{inc}}+\E_{\mathrm{sc}},
    \label{eq:incident-total}
\end{equation}
but \(\E_{\mathrm{inc}}\) is only the prescribed external drive.  The
induced field \(\E_{\mathrm{sc}}\), including the field created by the induced
charge, is part of the field that drives the electron fluid.

The local Drude model assumes
\begin{equation}
    \J=\sigma_D(\omega)\E,
    \qquad
    \sigma_D(\omega)=\frac{\ii\epsilon_0\omegap^2}{\omega+\ii\gamma},
    \label{eq:local-drude-current}
\end{equation}
where \(\gamma\) is a phenomenological momentum-relaxation rate and
\(\omegap\) is the plasma frequency.  It contains no spatial derivative of
the electron density and therefore no intrinsic microscopic length scale.

The hydrodynamic model adds the pressure force of a compressible electron
fluid.  In the linearized frequency-domain form its central equation is
\begin{equation}
    \boxed{
    \beta^2\nabla\!\left(\nabla\!\cdot\!\J\right)
    +\omega(\omega+\ii\gamma)\J
    =\ii\omega\epsilon_0\omegap^2\E.
    }
    \label{eq:hdm-current-central}
\end{equation}
The first term is the nonlocal correction.  It acts only on the
longitudinal part of the current, because it contains \(\nabla\!\cdot\!\J\).

The adjective hard-wall refers to the additional condition
\begin{equation}
    \boxed{\nvec\!\cdot\!\J=0\quad\text{at the metal boundary}.}
    \label{eq:hard-wall-central}
\end{equation}
It says that conduction electrons cannot flow through the geometrical
surface.  It does \emph{not} say that \(\nvec\!\cdot\!\E=0\).

\section{Why introduce a hydrodynamic model?}

\subsection{What local Drude theory can already do}

Local Drude theory can describe collective charge oscillations at an interface.
For example, a planar metal--dielectric interface supports a surface plasmon
whose charge is represented, in the local idealization, by an oscillating
surface sheet.  A nanoparticle can similarly support a dipolar plasmon, with
positive and negative induced charge patches on opposite sides.

The restoring force of these modes is the Coulomb field created by the charge
separation.  Fermi pressure is not required for the existence of the surface
plasmon.

\subsection{What is missing from local Drude theory}

Local Drude theory has two related limitations:
\begin{enumerate}
    \item the induced charge is compressed into an ideal mathematical sheet at
    an interface; and
    \item the formal bulk condition \(\eps_D(\omega)=0\) gives a plasma
    frequency but no finite wave number or spatial structure for the bulk
    longitudinal oscillation.
\end{enumerate}

The hydrodynamic pressure term supplies the missing spatial physics.  A density
gradient costs energy, so charge cannot be compressed arbitrarily sharply.
The induced surface charge becomes a narrow but finite-width layer, and bulk
longitudinal plasmons acquire the dispersion
\begin{equation}
    \omega(\omega+\ii\gamma)
    =\frac{\omegap^2}{\epsinf}+\beta^2 k^2.
    \label{eq:bulk-dispersion-preview}
\end{equation}

Vanilla HDM is therefore a mesoscopic extension of local Drude theory: it
retains the classical electromagnetic field and adds one coarse-grained
quantum ingredient, namely the compressibility of the conduction-electron gas.
It does not include equilibrium-density spill-out, tunnelling, or a fully
self-consistent exchange--correlation functional.

\section{Starting point: a charged compressible electron fluid}

\subsection{Fluid variables}

Let \(n(\mathbf r,t)\) be the conduction-electron number density and
\(\mathbf v(\mathbf r,t)\) their velocity.  The electron charge is \(-e\),
where \(e>0\).  The free-electron charge and current densities are
\begin{equation}
    \rho_f=-e n,
    \qquad
    \J=-e n\mathbf v.
    \label{eq:rho-J-definitions}
\end{equation}

The equilibrium jellium density is taken to be constant inside the metal,
\(n=n_0\), and zero outside.  In the vanilla hard-wall model this sharp
geometrical cutoff is imposed as an idealization:
\begin{equation}
    n_0(\mathbf r)=
    \begin{cases}
       n_0, & \mathbf r\in\Omega_m,\\
       0, & \mathbf r\notin\Omega_m.
    \end{cases}
    \label{eq:sharp-equilibrium-density}
\end{equation}
There is no equilibrium spill-out in this model.

\subsection{Continuity equation}

Conservation of the number of conduction electrons gives
\begin{equation}
    \partial_t n+\nabla\!\cdot(n\mathbf v)=0.
    \label{eq:continuity-n}
\end{equation}
Using \(\J=-en\mathbf v\), the same equation becomes charge continuity:
\begin{equation}
    \partial_t\rho_f+\nabla\!\cdot\!\J=0.
    \label{eq:continuity-rho}
\end{equation}
For harmonic fields,
\begin{equation}
    -\ii\omega\rho_f+\nabla\!\cdot\!\J=0,
    \qquad
    \rho_f=\frac{1}{\ii\omega}\nabla\!\cdot\!\J.
    \label{eq:rho-from-J}
\end{equation}

This equation is the first reason the divergence of the current is central:
charge is created wherever the current converges or terminates.  A current
may be large without creating charge if it is divergence-free.

\subsection{Momentum equation and pressure force}

The damped Euler equation for the electron fluid is
\begin{equation}
    m n\left(\partial_t+\mathbf v\!\cdot\!\nabla\right)\mathbf v
    =-en\E-m\gamma n\mathbf v-\nabla p(n),
    \label{eq:euler-full}
\end{equation}
where \(p(n)\) is the thermodynamic pressure of the electron gas.  The three
terms on the right are, respectively:
\begin{itemize}
    \item the electric force on the negatively charged electrons;
    \item phenomenological momentum relaxation; and
    \item the pressure-gradient force, which opposes compression.
\end{itemize}

We linearize around a uniform, stationary equilibrium:
\begin{equation}
    n=n_0+n_1,
    \qquad
    \mathbf v=\mathbf v_1,
    \qquad
    |n_1|\ll n_0.
    \label{eq:linearization}
\end{equation}
The convective term \(\mathbf v\!\cdot\!\nabla\mathbf v\) is second order and
is dropped.  The pressure is expanded as
\begin{equation}
    p(n_0+n_1)=p(n_0)+
    \left(\frac{\partial p}{\partial n}\right)_0n_1+O(n_1^2).
    \label{eq:pressure-linearization}
\end{equation}
Thus the linearized equations are
\begin{align}
    \partial_t n_1+n_0\nabla\!\cdot\mathbf v_1&=0,
    \label{eq:linear-continuity}\\[2mm]
    \partial_t\mathbf v_1+\gamma\mathbf v_1
    &=-\frac{e}{m}\E
    -\frac{1}{m n_0}
      \left(\frac{\partial p}{\partial n}\right)_0
      \nabla n_1.
    \label{eq:linear-euler}
\end{align}

Define the hydrodynamic velocity parameter \(\beta\) by
\begin{equation}
    \boxed{
    \beta^2\equiv
    \frac{1}{m}\left(\frac{\partial p}{\partial n}\right)_0.
    }
    \label{eq:beta-thermodynamic}
\end{equation}
Then the pressure force per electron mass is
\begin{equation}
    -\frac{\beta^2}{n_0}\nabla n_1.
    \label{eq:pressure-force-beta}
\end{equation}

Physically, \(\beta\) is not the drift velocity of the electron gas.  It is
the velocity scale associated with the restoring force generated by a density
gradient.  It plays the role of a sound velocity for the electron fluid, but
the relevant sound is a collective electronic density oscillation rather
than an ordinary acoustic wave.

\section{Deriving the hydrodynamic current equation}

\subsection{From velocity to current}

Inside the uniform metal, the linearized current is
\begin{equation}
    \J=-e n_0\mathbf v_1.
    \label{eq:linear-current-velocity}
\end{equation}
Multiply Eq.~\eqref{eq:linear-euler} by \(-e n_0\).  The left-hand side
becomes
\begin{equation}
    -e n_0(\partial_t+\gamma)\mathbf v_1
    =(\partial_t+\gamma)\J.
    \label{eq:current-time-left}
\end{equation}
The electric force becomes
\begin{equation}
    (-e n_0)\left(-\frac{e}{m}\E\right)
    =\frac{e^2n_0}{m}\E.
    \label{eq:current-electric-force}
\end{equation}
Introduce the plasma frequency
\begin{equation}
    \omegap^2\equiv\frac{e^2n_0}{m\epsilon_0}.
    \label{eq:plasma-frequency}
\end{equation}
Then the electric term is \(\epsilon_0\omegap^2\E\).

For the pressure term, use \(\J=-en_0\mathbf v_1\) in the linearized
continuity equation:
\begin{align}
    \partial_t n_1+n_0\nabla\!\cdot\mathbf v_1&=0,\\
    \partial_t n_1-\frac{1}{e}\nabla\!\cdot\!\J&=0.
    \label{eq:n1-continuity-J}
\end{align}
Therefore,
\begin{equation}
    \partial_t n_1=\frac{1}{e}\nabla\!\cdot\!\J.
    \label{eq:n1-time-J}
\end{equation}
The pressure contribution after multiplication by \(-e n_0\) is
\begin{equation}
    (-e n_0)\left(-\frac{\beta^2}{n_0}\nabla n_1\right)
    =e\beta^2\nabla n_1.
    \label{eq:pressure-current-term}
\end{equation}

In the frequency domain, Eq.~\eqref{eq:n1-continuity-J} gives
\begin{equation}
    -\ii\omega n_1=\frac{1}{e}\nabla\!\cdot\!\J,
    \qquad
    n_1=\frac{\ii}{e\omega}\nabla\!\cdot\!\J.
    \label{eq:n1-frequency}
\end{equation}
Consequently,
\begin{equation}
    e\beta^2\nabla n_1
    =\frac{\ii\beta^2}{\omega}
      \nabla\!\left(\nabla\!\cdot\!\J\right).
    \label{eq:pressure-frequency-current}
\end{equation}

The frequency-domain current equation is therefore
\begin{equation}
    (\gamma-\ii\omega)\J
    =\epsilon_0\omegap^2\E
    +\frac{\ii\beta^2}{\omega}
      \nabla\!\left(\nabla\!\cdot\!\J\right).
    \label{eq:current-before-final}
\end{equation}
Multiplying by \(\ii\omega\) gives
\begin{equation}
    \boxed{
    \beta^2\nabla\!\left(\nabla\!\cdot\!\J\right)
    +\omega(\omega+\ii\gamma)\J
    =\ii\omega\epsilon_0\omegap^2\E.
    }
    \label{eq:current-final-derived}
\end{equation}

The sign of the pressure term is worth interpreting.  In Fourier space,
\(\nabla\to\ii\mathbf k\), so
\begin{equation}
    \nabla(\nabla\!\cdot\!\J)
    \longrightarrow
    -\mathbf k(\mathbf k\!\cdot\!\J).
    \label{eq:grad-div-fourier}
\end{equation}
For a longitudinal current, this is \(-k^2\J_L\).  The pressure term thus
opposes short-wavelength compression and changes the longitudinal resonance
condition by a positive term \(\beta^2 k^2\) in the dispersion relation.

\subsection{Local Drude limit and polarization form}

If \(\beta=0\), Eq.~\eqref{eq:current-final-derived} becomes
\begin{equation}
    \J=\frac{\ii\epsilon_0\omegap^2}{\omega+\ii\gamma}\E
    \equiv\sigma_D(\omega)\E,
    \label{eq:drude-from-hdm}
\end{equation}
which is precisely the local Drude current.

The free-electron polarization is defined by
\begin{equation}
    \J=\partial_t\Pfield_f=-\ii\omega\Pfield_f,
    \qquad
    \Pfield_f=\frac{\J}{-\ii\omega}.
    \label{eq:free-polarization-current}
\end{equation}
The total displacement inside the metal is split as
\begin{equation}
    \Dfield_m
    =\epsilon_0\epsinf\E+\Pfield_f.
    \label{eq:D-split}
\end{equation}
The first term is the local background response; the second is the response
of the hydrodynamic conduction electrons.

In the local limit, substituting Eq.~\eqref{eq:drude-from-hdm} into
Eq.~\eqref{eq:free-polarization-current} gives
\begin{equation}
    \Pfield_f
    =-\epsilon_0
      \frac{\omegap^2}{\omega(\omega+\ii\gamma)}\E.
    \label{eq:local-free-polarization}
\end{equation}
Therefore the familiar local Drude permittivity is
\begin{equation}
    \boxed{
    \eps_D(\omega)
    =\epsinf-
      \frac{\omegap^2}{\omega(\omega+\ii\gamma)}.
    }
    \label{eq:drude-permittivity}
\end{equation}

\section{What should \texorpdfstring{$\beta$}{beta} be?}

The hydrodynamic equation itself does not determine \(\beta\).  It is a
long-wavelength closure of a microscopic electron response, so its value
depends on which limit of the electron gas one is trying to reproduce.

\subsection{Static compressibility: \texorpdfstring{$\beta^2=\vF^2/3$}{beta2=vF2/3}}

For a three-dimensional ideal Fermi gas at zero temperature,
\begin{equation}
    n=\frac{k_F^3}{3\pi^2},
    \qquad
    \vF=\frac{\hbar k_F}{m},
    \qquad
    E_F=\frac{1}{2}m\vF^2.
    \label{eq:fermi-definitions}
\end{equation}
Its pressure is
\begin{equation}
    p=\frac{2}{5}nE_F.
    \label{eq:fermi-pressure}
\end{equation}
Since \(E_F\propto n^{2/3}\), differentiation gives
\begin{align}
    \frac{\partial p}{\partial n}
    &=\frac{2}{3}E_F
      =\frac{1}{3}m\vF^2,\\
    \beta_{\mathrm{static}}^2
    &=\frac{1}{m}\frac{\partial p}{\partial n}
      =\frac{\vF^2}{3}.
    \label{eq:beta-static}
\end{align}

This is the equilibrium compressibility: it asks how much pressure is produced
by a slowly varying change in density after the electron gas has time to
re-equilibrate locally.

\subsection{Collisionless optical response: \texorpdfstring{$\beta^2=3\vF^2/5$}{beta2=3vF2/5}}

At optical frequencies, the perturbation can be collisionless on the time
scale of the oscillation.  The electron distribution then responds through
Fermi-surface kinematics rather than through static thermodynamic compression.
The long-wavelength expansion of the collisionless Lindhard response has the
form
\begin{equation}
    \eps_L(k,\omega)-1
    \simeq
    -\frac{\omegap^2}{\omega^2}
    \left[1+\frac{3}{5}\frac{k^2\vF^2}{\omega^2}+\cdots\right].
    \label{eq:lindhard-expansion}
\end{equation}
The corresponding bulk-plasmon dispersion is
\begin{equation}
    \omega^2
    =\omegap^2+\frac{3}{5}\vF^2 k^2+\cdots
    \qquad (\epsinf=1,\ \gamma=0).
    \label{eq:collisionless-dispersion}
\end{equation}
Matching the hydrodynamic dispersion to this expansion gives
\begin{equation}
    \boxed{\beta_{\mathrm{collisionless}}^2=\frac{3}{5}\vF^2.}
    \label{eq:beta-collisionless}
\end{equation}

The factors \(1/3\) and \(3/5\) are therefore not contradictory estimates of
one universal number.  They describe different limits:
\begin{equation}
    \frac{1}{3}\vF^2
    \quad\text{comes from static compressibility},
    \qquad
    \frac{3}{5}\vF^2
    \quad\text{from collisionless dynamic response}.
    \label{eq:beta-two-limits}
\end{equation}

For optical plasmonics, the conventional vanilla HDM choice is often
\(\beta^2=3\vF^2/5\), because the relevant response is usually closer to the
collisionless regime.  The static value is also used when the model is meant
to represent a quasi-static compressibility.

\subsection{A frequency-dependent interpolation}

A more refined long-wavelength matching to a conserving relaxation-time
response gives the Halevi interpolation \cite{Halevi1995,Mermin1970}.  With
the present convention and a relaxation rate \(\gamma\), it can be written
as
\begin{equation}
    \boxed{
    \beta^2(\omega)
    =\vF^2
      \frac{\frac{3}{5}\omega+\frac{\ii\gamma}{3}}
           {\omega+\ii\gamma}.
    }
    \label{eq:halevi-beta}
\end{equation}
Indeed,
\begin{align}
    \omega\gg\gamma&:\quad
    \beta^2(\omega)\longrightarrow\frac{3}{5}\vF^2,\\
    \omega\ll\gamma&:\quad
    \beta^2(\omega)\longrightarrow\frac{1}{3}\vF^2.
    \label{eq:halevi-limits}
\end{align}

The interpolation is not an arbitrary linear blend.  It is obtained by
matching the order-\(k^2\) expansion of the hydrodynamic response to a kinetic
relaxation-time response with the Mermin correction, which preserves local
particle-number conservation.  At intermediate \(\omega/\gamma\), \(\beta^2\)
is complex.  This means that the pressure-like part of the response is not
exactly in phase with the density; collisions produce a phase lag.

In practice there are two separate approximations here:
\begin{enumerate}
    \item a long-wavelength expansion in \(k\); and
    \item a frequency interpolation in \(\omega/\gamma\).
\end{enumerate}
The standard vanilla implementation usually takes \(\beta\) to be a real
constant, most commonly \(\sqrt{3/5}\,\vF\).  Using Eq.~\eqref{eq:halevi-beta}
makes the hydrodynamic coefficient frequency-dependent and possibly complex,
which is allowed but is a more refined constitutive model.

\section{Transverse and longitudinal response in a homogeneous metal}

\subsection{The decomposition is relative to a spatial wavevector}

Consider a homogeneous region and a Fourier component with wavevector
\(\mathbf k\).  Any vector field can be decomposed as
\begin{equation}
    \E=\E_T+\E_L,
    \qquad
    \J=\J_T+\J_L,
    \label{eq:T-L-decomposition}
\end{equation}
where
\begin{equation}
    \mathbf k\!\cdot\!\E_T=0,
    \quad
    \mathbf k\times\E_L=0,
    \qquad
    \mathbf k\!\cdot\!\J_T=0,
    \quad
    \mathbf k\times\J_L=0.
    \label{eq:T-L-definitions}
\end{equation}
Thus longitudinal means parallel to the relevant spatial wavevector.  It
does not mean parallel to the propagation direction of an externally incident
electromagnetic wave unless those directions happen to coincide.

In a finite nanoparticle there need not be one unique wavevector.  The
transverse/longitudinal language then refers to the Helmholtz decomposition or
to the modal Fourier components of the fields.  For a planar interface, the
in-plane wavevector is conserved, but the transverse and longitudinal modes
usually have different normal wavevectors.

\subsection{Transverse current: unchanged Drude response}

In Fourier space,
\begin{equation}
    \nabla(\nabla\!\cdot\!\J)
    \longrightarrow -\mathbf k(\mathbf k\!\cdot\!\J).
    \label{eq:pressure-fourier-again}
\end{equation}
For the transverse part, \(\mathbf k\!\cdot\!\J_T=0\), so the pressure term
vanishes.  Equation~\eqref{eq:current-final-derived} becomes
\begin{equation}
    \J_T
    =\frac{\ii\epsilon_0\omegap^2}{\omega+\ii\gamma}\E_T.
    \label{eq:JT}
\end{equation}
Therefore
\begin{equation}
    \boxed{\epst(\omega)=\eps_D(\omega).}
    \label{eq:epsilon-transverse}
\end{equation}
The hydrodynamic correction does not modify a divergence-free transverse
current in a homogeneous bulk.

\subsection{Longitudinal current: nonlocal response}

For the longitudinal part,
\(\nabla(\nabla\!\cdot\!\J_L)\to-k^2\J_L\).  The current equation gives
\begin{equation}
    \left[\omega(\omega+\ii\gamma)-\beta^2k^2\right]\J_L
    =\ii\omega\epsilon_0\omegap^2\E_L.
    \label{eq:JL-equation}
\end{equation}
Hence
\begin{equation}
    \J_L
    =\frac{\ii\omega\epsilon_0\omegap^2}
           {\omega(\omega+\ii\gamma)-\beta^2k^2}\E_L.
    \label{eq:JL-response}
\end{equation}
Using \(\Dfield=\epsilon_0\epsinf\E+\Pfield_f\) and
\(\Pfield_f=\J/(-\ii\omega)\), the longitudinal dielectric function is
\begin{equation}
    \boxed{
    \epsl(k,\omega)
    =\epsinf-
      \frac{\omegap^2}
      {\omega(\omega+\ii\gamma)-\beta^2k^2}.
    }
    \label{eq:epsilon-longitudinal}
\end{equation}
This is spatial dispersion: the response depends on both frequency and
wavevector.

\subsection{Bulk longitudinal plasmons}

A longitudinal bulk mode has \(\Hfield_L=0\) and, in the absence of an
external charge, satisfies
\begin{equation}
    \nabla\!\cdot\!\Dfield=0.
    \label{eq:longitudinal-gauss}
\end{equation}
For a nonzero longitudinal electric field this requires
\(\epsl(k,\omega)=0\).  Therefore
\begin{equation}
    \boxed{
    \omega(\omega+\ii\gamma)
    =\frac{\omegap^2}{\epsinf}+\beta^2k^2.
    }
    \label{eq:bulk-hdm-dispersion}
\end{equation}
Without damping,
\begin{equation}
    \omega_L(k)=
    \sqrt{\frac{\omegap^2}{\epsinf}+\beta^2k^2}.
    \label{eq:bulk-hdm-dispersion-lossless}
\end{equation}

This is the bulk plasmon branch supplied by HDM.  Local Drude theory has the
formal condition \(\eps_D(\omega)=0\), but because it contains no \(k^2\) term
it cannot determine a longitudinal wavelength or spatial profile.

\section{Incident field, total field, and the origin of the longitudinal drive}

It is useful to separate three ideas that are often accidentally conflated.

\subsection{Incident electromagnetic wave}

For a homogeneous incident plane wave,
\begin{equation}
    \E_{\mathrm{inc}}(\mathbf r)
    =\E_0\ee^{\ii\mathbf k_{\mathrm{inc}}\cdot\mathbf r},
    \qquad
    \mathbf k_{\mathrm{inc}}\!\cdot\!\E_0=0.
    \label{eq:incident-plane-wave}
\end{equation}
It is transverse relative to its own propagation wavevector.  In homogeneous
bulk, the induced local Drude current is also transverse and therefore has
\(\nabla\!\cdot\!\J=0\).  There is no bulk charge-density wave at that same
plane-wave wavevector.

At normal incidence on an infinite planar interface, take
\(\mathbf k_{\mathrm{inc}}=k\hat{\mathbf z}\) and tangential electric field
\(\E_0\parallel\) interface.  Even if the tangential current varies with
depth because of attenuation,
\begin{equation}
    \nabla\!\cdot\!\J
    =\partial_zJ_z+\nabla_\parallel\!\cdot\!\J_\parallel=0
    \label{eq:normal-incidence-no-charge}
\end{equation}
for the ideal translationally invariant problem: \(J_z=0\) and the tangential
current has no in-plane divergence.  Energy can still be dissipated through
\(\gamma\); charge compression is simply not required.

\subsection{Oblique incidence and boundary launching}

For oblique TM incidence at a planar interface,
\begin{equation}
    \mathbf k_{\mathrm{inc}}=(q,0,k_{z,\mathrm{inc}}),
    \qquad
    \E_0=(E_x,0,E_z),
    \qquad
    qE_x+k_{z,\mathrm{inc}}E_z=0.
    \label{eq:TM-incident}
\end{equation}
The normal field component drives a normal current.  The hard wall prevents
that current from crossing the interface:
\begin{equation}
    J_{z,T}(0)+J_{z,L}(0)=0.
    \label{eq:current-cancellation-wall}
\end{equation}
The longitudinal component is therefore generated self-consistently so that
the total normal electron flux vanishes.  This is the precise sense in which
the boundary launches the longitudinal HDM response.

The longitudinal mode has
\begin{equation}
    \mathbf k_L=(q,0,k_{z,L}),
    \qquad
    \E_L\parallel\mathbf k_L,
    \label{eq:planar-longitudinal-wavevector}
\end{equation}
but \(k_{z,L}\) is determined by the longitudinal dispersion
Eq.~\eqref{eq:bulk-hdm-dispersion}, not by the incident electromagnetic
dispersion.  Only the tangential momentum \(q\) is shared by the fields due to
translational invariance along the interface.  Consequently,
\begin{equation}
    \boxed{\E_0\parallel\mathbf k_L\text{ is not a general relation}.}
    \label{eq:no-general-parallel}
\end{equation}

The longitudinal field is not a second externally incident electromagnetic
wave.  It is part of \(\E_{\mathrm{tot}}\), produced by the induced charge and
required by the boundary condition.

\section{Surface charge in local Drude theory and finite-width charge in HDM}

\subsection{Why local Drude theory has surface charge}

Let a planar metal occupy \(z<0\).  In local Drude theory the current can be
written schematically as
\begin{equation}
    \J(\mathbf r)=\Theta(-z)\J_m(\mathbf r),
    \label{eq:step-current}
\end{equation}
where \(\Theta\) is the step function.  Its divergence is
\begin{align}
    \nabla\!\cdot\!\J
    &=\nabla\!\cdot\!\left[\Theta(-z)\J_m\right]\\
    &=\Theta(-z)\nabla\!\cdot\!\J_m
      -J_{m,z}(0)\delta(z).
    \label{eq:step-current-divergence}
\end{align}
The delta function is the mathematical representation of an induced surface
charge:
\begin{equation}
    \rho_f=\rho_{\mathrm{vol}}+\sigma_s(x,y)\delta(z),
    \qquad
    \sigma_s=-\frac{J_{m,z}(0)}{\ii\omega}.
    \label{eq:local-surface-charge}
\end{equation}

Thus surface charge does not mean that local Drude theory violates charge
continuity.  The charge appears because the current is terminated at the
geometrical boundary.  Local Drude has no pressure-gradient energy and no
microscopic length with which to resolve the termination, so the charge is
represented as an infinitely thin sheet.

\subsection{Local planar surface plasmon}

In the quasistatic limit, take potentials of the form
\begin{equation}
    \phi_m=A\ee^{\ii qx+qz},\quad z<0,
    \qquad
    \phi_d=A\ee^{\ii qx-qz},\quad z>0.
    \label{eq:surface-potentials}
\end{equation}
Both potentials satisfy Laplace's equation away from the interface, so there
is no volume charge in either homogeneous region.  Nevertheless, the fields
are normal to the interface with opposite signs:
\begin{equation}
    E_{m,z}(0)=-qA,
    \qquad
    E_{d,z}(0)=+qA.
    \label{eq:surface-normal-fields}
\end{equation}
Continuity of the normal displacement gives
\begin{equation}
    \eps_D(\omega)E_{m,z}=\epsd E_{d,z},
    \qquad
    \boxed{\eps_D(\omega)+\epsd=0.}
    \label{eq:local-surface-plasmon-condition}
\end{equation}
For a lossless Drude metal,
\begin{equation}
    \eps_D(\omega)=\epsinf-\frac{\omegap^2}{\omega^2},
    \qquad
    \omega_{\mathrm{sp}}
    =\frac{\omegap}{\sqrt{\epsinf+\epsd}}.
    \label{eq:surface-plasmon-frequency}
\end{equation}

The surface plasmon is a self-consistent oscillation: the charge produces an
electric field, the field drives electron motion, and the resulting current
recreates the charge.  The Coulomb restoring force is sufficient for this
mode; HDM is not needed for its existence.

\subsection{What HDM changes}

In HDM the pressure force penalizes a sharp density gradient.  The charge that
was a delta-function sheet in local Drude theory becomes a narrow dynamical
layer inside the metal.  To see its characteristic scale, take the divergence
of Eq.~\eqref{eq:current-final-derived}:
\begin{equation}
    \beta^2\nabla^2(\nabla\!\cdot\!\J)
    +\omega(\omega+\ii\gamma)\nabla\!\cdot\!\J
    =\ii\omega\epsilon_0\omegap^2\nabla\!\cdot\!\E.
    \label{eq:div-current-equation}
\end{equation}
Using \(\nabla\!\cdot\!\J=\ii\omega\rho_f\), we need a relation between
\(\nabla\!\cdot\!\E\) and \(\rho_f\).  In the absence of externally imposed
free charge, Ampere's law implies
\begin{equation}
    \nabla\!\cdot\!\J
    =\ii\omega\epsilon_0\epsinf\nabla\!\cdot\!\E,
    \label{eq:divJ-divE}
\end{equation}
and therefore
\begin{equation}
    \rho_f=\epsilon_0\epsinf\nabla\!\cdot\!\E.
    \label{eq:rho-divE}
\end{equation}
Substitution into Eq.~\eqref{eq:div-current-equation} yields
\begin{equation}
    \boxed{
    \left[
      \beta^2\nabla^2
      +\omega(\omega+\ii\gamma)
      -\frac{\omegap^2}{\epsinf}
    \right]\rho_f=0.
    }
    \label{eq:rho-wave-equation}
\end{equation}
Equivalently,
\begin{equation}
    \left(\nabla^2+k_L^2\right)\rho_f=0,
    \qquad
    k_L^2=
    \frac{\omega(\omega+\ii\gamma)-\omegap^2/\epsinf}{\beta^2}.
    \label{eq:rho-kL}
\end{equation}

For frequencies below the bulk longitudinal-plasmon frequency, \(k_L^2\) is
typically negative.  Writing \(k_L=\ii\kappa_L\), the density decays into the
metal over
\begin{equation}
    \ell_L\sim\kappa_L^{-1},
    \qquad
    \kappa_L^2\simeq
    \frac{\omegap^2/\epsinf-\omega(\omega+\ii\gamma)}{\beta^2}.
    \label{eq:charge-layer-length}
\end{equation}
The rough scale \(\beta/\omegap\) is often quoted, although the precise decay
length depends on frequency, \(\epsinf\), damping, and geometry.

The distinction can be summarized as follows.
\begin{center}
\begin{tabular}{>{\raggedright\arraybackslash}p{0.43\textwidth}
                >{\raggedright\arraybackslash}p{0.43\textwidth}}
\toprule
\textbf{Local Drude theory} & \textbf{Hard-wall HDM}\\
\midrule
Surface charge represented by \(\sigma_s\delta(s)\) &
Induced charge occupies a finite layer inside the metal\\
No pressure length scale &
Pressure introduces the scale \(\ell_L\sim\beta/\omegap\)\\
Surface plasmons exist through Coulomb restoring force &
Surface plasmons are shifted and coupled to longitudinal response\\
Formal bulk-plasmon frequency but no finite-\(k\) branch &
Bulk longitudinal plasmons disperse as \(\omega^2\sim\omegap^2/\epsinf+\beta^2k^2\)\\
\bottomrule
\end{tabular}
\end{center}

The hard-wall charge layer is confined to the metal side of the geometrical
surface.  This is different from self-consistent hydrodynamic or quantum
hydrodynamic models, where the equilibrium density can spill out and the
induced charge can extend into the dielectric region.

\section{Eliminating \texorpdfstring{$\J$}{J}: an electric-field-only equation}

Many numerical implementations solve directly for \(\E\).  We now derive the
corresponding equation carefully, because the term
\(\nabla(\nabla\!\cdot\!\E)\) comes specifically from eliminating \(\J\) with
Maxwell's equations.

\subsection{Maxwell equations inside the metal}

With the convention \eqref{eq:time-convention}, Maxwell's equations are
\begin{align}
    \nabla\!\times\!\E&=\ii\omega\mu_0\Hfield,
    \label{eq:faraday-frequency}\\
    \nabla\!\times\!\Hfield
    &=\J-\ii\omega\epsilon_0\epsinf\E.
    \label{eq:ampere-frequency}
\end{align}
The second equation contains the conduction current separately because the
background polarization has already been included in \(\epsinf\).

Use Faraday's law to write
\begin{equation}
    \Hfield=\frac{1}{\ii\omega\mu_0}\nabla\!\times\!\E.
    \label{eq:H-from-E}
\end{equation}
Substitute this into Ampere's law:
\begin{align}
    \frac{1}{\ii\omega\mu_0}
    \nabla\!\times\!\nabla\!\times\!\E
    &=\J-\ii\omega\epsilon_0\epsinf\E,\\
    \J
    &=\frac{1}{\ii\omega\mu_0}
      \left[
        \nabla\!\times\!\nabla\!\times\!\E
        -\epsinf k_0^2\E
      \right],
    \qquad
    k_0=\frac{\omega}{c}.
    \label{eq:J-from-E}
\end{align}
This is the exact Maxwell relation used to remove \(\J\).

Now take the divergence.  Since the divergence of any curl is zero,
\begin{align}
    \nabla\!\cdot\!\J
    &=\frac{1}{\ii\omega\mu_0}
      \left[
        \underbrace{\nabla\!\cdot\!
        (\nabla\!\times\!\nabla\!\times\!\E)}_{0}
        -\epsinf k_0^2\nabla\!\cdot\!\E
      \right]\\
    &=-\frac{\epsinf k_0^2}{\ii\omega\mu_0}
      \nabla\!\cdot\!\E
     =\ii\omega\epsilon_0\epsinf
      \nabla\!\cdot\!\E.
    \label{eq:divJ-from-E}
\end{align}
The last equality uses \(k_0^2=\omega^2\mu_0\epsilon_0\).

\subsection{Substitution into the hydrodynamic equation}

Divide the current equation \eqref{eq:current-final-derived} by
\(\omega(\omega+\ii\gamma)\).  Defining
\begin{equation}
    a\equiv\frac{\beta^2}{\omega(\omega+\ii\gamma)},
    \label{eq:a-definition}
\end{equation}
we have
\begin{equation}
    a\nabla(\nabla\!\cdot\!\J)+\J=\sigma_D\E.
    \label{eq:current-a-form}
\end{equation}
Using Eq.~\eqref{eq:divJ-from-E}, the current can equivalently be written
\begin{align}
    \J
    &=\sigma_D\E-a\nabla(\nabla\!\cdot\!\J)\\
    &=\sigma_D\E
      -\frac{\ii\epsilon_0\epsinf\beta^2}{\omega+\ii\gamma}
       \nabla(\nabla\!\cdot\!\E).
    \label{eq:J-electric-constitutive}
\end{align}
This equation already shows that the correction is purely longitudinal: if
\(\nabla\!\cdot\!\E=0\), it reduces to local Drude response.

Substitute Eq.~\eqref{eq:J-from-E} and Eq.~\eqref{eq:divJ-from-E} directly
into Eq.~\eqref{eq:current-a-form}:
\begin{align}
    &\frac{1}{\ii\omega\mu_0}
      \left[\nabla\!\times\!\nabla\!\times\!\E
      -\epsinf k_0^2\E\right]
    -a\frac{\epsinf k_0^2}{\ii\omega\mu_0}
      \nabla(\nabla\!\cdot\!\E)
    =\sigma_D\E.
    \label{eq:substitution-before-multiply}
\end{align}
Multiplying by \(\ii\omega\mu_0\) gives
\begin{equation}
    \nabla\!\times\!\nabla\!\times\!\E
    -\epsinf k_0^2\E
    -a\epsinf k_0^2\nabla(\nabla\!\cdot\!\E)
    =\ii\omega\mu_0\sigma_D\E.
    \label{eq:E-equation-intermediate}
\end{equation}
Now
\begin{equation}
    \ii\omega\mu_0\sigma_D
    =-k_0^2\frac{\omegap^2}{\omega(\omega+\ii\gamma)}.
    \label{eq:sigma-conversion}
\end{equation}
Combining the two local terms produces \(\eps_D\), and we obtain
\begin{equation}
    \boxed{
    \nabla\!\times\!\nabla\!\times\!\E
    -k_0^2\eps_D(\omega)\E
    =\frac{\epsinf k_0^2\beta^2}
           {\omega(\omega+\ii\gamma)}
      \nabla(\nabla\!\cdot\!\E).
    }
    \label{eq:E-only-HDM}
\end{equation}

The term on the right is absent in local electrodynamics.  It is precisely the
electric-field representation of the pressure response.  Since
\(\nabla\!\cdot\!\E_T=0\), it acts only on \(\E_L\).  The field in this
equation is the total field, not just the incident field.

\section{Boundary conditions for the hard-wall model}

The hydrodynamic equation is a higher-order spatial equation than local Drude
theory, so Maxwell's usual interface conditions are not sufficient by
themselves.  The electron fluid supplies one additional boundary condition.

\subsection{Maxwell boundary conditions}

Assume no externally imposed singular surface current or external free surface
charge.  The integral Maxwell laws give
\begin{align}
    \nvec\times(\E_d-\E_m)&=0,
    \label{eq:bc-E-tangential}\\
    \nvec\times(\Hfield_d-\Hfield_m)&=0,
    \label{eq:bc-H-tangential}\\
    \nvec\!\cdot(\Bfield_d-\Bfield_m)&=0,
    \label{eq:bc-B-normal}\\
    \nvec\!\cdot(\Dfield_d-\Dfield_m)&=0.
    \label{eq:bc-D-normal}
\end{align}
The last equation is written for the total displacement, including the
hydrodynamic free-electron polarization inside the metal.

Outside the metal,
\begin{equation}
    \Dfield_d=\epsilon_0\epsd\E_d.
    \label{eq:D-outside}
\end{equation}
Inside,
\begin{equation}
    \Dfield_m=\epsilon_0\epsinf\E_m+\Pfield_f,
    \qquad
    \Pfield_f=\frac{\J}{-\ii\omega}.
    \label{eq:D-inside-again}
\end{equation}

\subsection{Derivation of the hard-wall additional boundary condition}

The hard-wall assumption is a statement about particle transport:
\begin{equation}
    \boxed{\nvec\!\cdot\!\J_m=0.}
    \label{eq:bc-hard-wall-J}
\end{equation}
Because \(\J=-\ii\omega\Pfield_f\), this is equivalent for \(\omega\neq0\)
to
\begin{equation}
    \boxed{\nvec\!\cdot\!\Pfield_f=0.}
    \label{eq:bc-hard-wall-P}
\end{equation}
Insert this into the normal-displacement condition:
\begin{align}
    \nvec\!\cdot\!\Dfield_d
    &=\nvec\!\cdot\!\Dfield_m,\\
    \epsilon_0\epsd E_{d,n}
    &=\epsilon_0\epsinf E_{m,n}+P_{f,n},\\
    \boxed{\epsd E_{d,n}=\epsinf E_{m,n}.}
    \label{eq:bc-normal-E-final}
\end{align}
For vacuum outside, \(\epsd=1\), so
\begin{equation}
    E_{d,n}=\epsinf E_{m,n}.
    \label{eq:vacuum-normal-E}
\end{equation}

The electric field inside the metal is therefore not multiplied by
\(\epsinf\).  Rather, the normal boundary relation contains \(\epsinf\)
because the normal component of the conduction-electron polarization vanishes
at a hard wall.  The remaining normal displacement is the background part
\(\epsilon_0\epsinf\E_m\).

The standard hard-wall interface conditions can consequently be listed as
\begin{equation}
    \boxed{
    \begin{aligned}
        \nvec\times(\E_d-\E_m)&=0,\\
        \nvec\times(\Hfield_d-\Hfield_m)&=0,\\
        \epsd E_{d,n}&=\epsinf E_{m,n},\\
        \nvec\!\cdot\!\J_m&=0.
    \end{aligned}}
    \label{eq:all-hard-wall-bcs}
\end{equation}

The last condition is the additional boundary condition (ABC) required by
HDM.  It is a no-flux condition, not a no-field condition:
\begin{equation}
    \nvec\!\cdot\!\J=0
    \quad\not\Rightarrow\quad
    \nvec\!\cdot\!\E=0.
    \label{eq:no-flux-not-no-field}
\end{equation}

\subsection{Electric-field-only form of the ABC}

When solving the electric-field equation \eqref{eq:E-only-HDM}, one can
replace \(\J\) using Eq.~\eqref{eq:J-from-E}.  The hard-wall condition becomes
\begin{equation}
    \boxed{
    \nvec\!\cdot
    \left[
      \nabla\!\times\!\nabla\!\times\!\E_m
      -\epsinf k_0^2\E_m
    \right]=0.
    }
    \label{eq:E-only-ABC}
\end{equation}
The omitted prefactor \(1/(\ii\omega\mu_0)\) is nonzero and therefore does
not affect the zero condition.

\subsection{Two bookkeeping conventions for surface charge}

There are two equivalent ways to describe interface charge, provided one does
not mix conventions.

In the total-displacement convention used above, the hydrodynamic conduction
electrons are included in \(\Dfield_m\), so there is no externally imposed
delta-function sheet in the normal \(\Dfield\) boundary condition.  Their
charge is encoded in \(\nabla\!\cdot\!\Pfield_f\) inside the metal.

Alternatively, one may treat \(\Pfield_f\) as a separate polarization and
write its divergence explicitly as the conduction-electron charge.  In local
Drude theory the sharp step in \(\Pfield_f\) produces a delta-function sheet.
In hard-wall HDM, \(P_{f,n}=0\) prevents such a singular normal-polarization
jump at the ideal wall, while the pressure term resolves the charge into a
finite-width volume layer just inside the metal.

\section{What exactly is \texorpdfstring{$\epsinf$}{epsilon-infinity}?}

\subsection{Physical meaning}

The decomposition
\begin{equation}
    \Dfield_m=\epsilon_0\epsinf\E+\Pfield_f
    \label{eq:epsinf-decomposition}
\end{equation}
separates two electronic responses:
\begin{enumerate}
    \item \(\Pfield_f\): the explicitly modeled mobile conduction-electron
    fluid; and
    \item \(\epsilon_0(\epsinf-1)\E\): a local background polarization due
    to bound/core electrons and interband transitions that are not included in
    the hydrodynamic fluid.
\end{enumerate}

For an ideal jellium gas in vacuum, there are no bound electrons in the model,
so
\begin{equation}
    \epsinf=1.
    \label{eq:jellium-epsinf}
\end{equation}
For a real metal, \(\epsinf\) is generally an effective material parameter.
It is not the full measured metal permittivity at the operating frequency.
The Drude part must not be counted twice.

\subsection{How \texorpdfstring{$\epsinf$}{epsilon-infinity} is obtained}

In the simplest Drude--HDM model, one fits the measured complex permittivity
over a selected frequency window to
\begin{equation}
    \eps_{\mathrm{exp}}(\omega)
    \simeq
    \epsinf-
    \frac{\omegap^2}{\omega(\omega+\ii\gamma)}.
    \label{eq:epsinf-fit}
\end{equation}
Depending on what is known independently, one may fit \(\epsinf\),
\(\omegap\), and \(\gamma\) simultaneously, or fix \(\omegap\) and
\(\gamma\) and fit only the background.

Equivalently, a provisional value can be extracted at a given frequency as
\begin{equation}
    \epsinf
    \simeq
    \eps_{\mathrm{exp}}(\omega)
    +\frac{\omegap^2}{\omega(\omega+\ii\gamma)},
    \label{eq:epsinf-extraction}
\end{equation}
although fitting over a frequency interval is more reliable.

If interband transitions are important, they should be modeled explicitly,
for example with a Drude--Lorentz form:
\begin{equation}
    \eps(\omega)
    =\epsinf-
    \frac{\omegap^2}{\omega(\omega+\ii\gamma)}
    +\sum_j\frac{f_j}
    {\omega_j^2-\omega^2-\ii\Gamma_j\omega}.
    \label{eq:drude-lorentz}
\end{equation}
In that case \(\epsinf\) is the residual high-frequency local background
after the explicitly included Drude and Lorentz contributions have been
separated.  In a fully realistic material model it can itself be complex and
frequency-dependent.

The practical rule is therefore:
\begin{equation}
    \boxed{\epsinf\text{ is fitted or calibrated, not a universal metal constant.}}
    \label{eq:epsinf-rule}
\end{equation}

\section{A compact physical picture}

The entire response mechanism can be read as a sequence:
\begin{enumerate}
    \item The total electric field drives the conduction-electron current.
    \item A divergence-free tangential current remains transverse and follows
    the local Drude response.
    \item A normal current driven toward a hard wall cannot cross it, so the
    current must terminate or be redirected near the boundary.
    \item Current convergence creates charge through
    \(\rho_f=(\ii\omega)^{-1}\nabla\!\cdot\!\J\).
    \item The charge creates a self-consistent longitudinal field
    \(\E_L=-\nabla\phi_L\).
    \item Coulomb forces produce the familiar surface-plasmon restoring force;
    finite pressure spreads the charge and adds the longitudinal \(\beta^2k^2\)
    dispersion.
    \item The hard-wall condition enforces zero normal electron flux while
    allowing a nonzero normal electric field.
\end{enumerate}

The two limiting pictures are therefore compatible:
\begin{equation}
    \boxed{
    \begin{gathered}
      \text{local Drude: ideal surface-charge sheet}\\[-1mm]
      \longrightarrow\\[-1mm]
      \text{HDM: finite-width dynamical charge layer}.
    \end{gathered}
    }
    \label{eq:final-picture}
\end{equation}

\section{The vanilla hard-wall HDM in one place}

For reference, the model consists of the following ingredients.

\subsection*{Bulk equations inside the metal}
\begin{align}
    \nabla\!\times\!\E&=\ii\omega\mu_0\Hfield,\\
    \nabla\!\times\!\Hfield
    &=\J-\ii\omega\epsilon_0\epsinf\E,\\
    \beta^2\nabla(\nabla\!\cdot\!\J)
    +\omega(\omega+\ii\gamma)\J
    &=\ii\omega\epsilon_0\omegap^2\E.
    \label{eq:model-bulk-summary}
\end{align}

\subsection*{Interface conditions}
\begin{align}
    \nvec\times(\E_d-\E_m)&=0,\\
    \nvec\times(\Hfield_d-\Hfield_m)&=0,\\
    \epsd E_{d,n}&=\epsinf E_{m,n},\\
    \nvec\!\cdot\!\J_m&=0.
    \label{eq:model-bc-summary}
\end{align}

\subsection*{Material parameters}
\begin{equation}
    \omegap^2=\frac{e^2n_0}{m\epsilon_0},
    \qquad
    \sigma_D=\frac{\ii\epsilon_0\omegap^2}{\omega+\ii\gamma},
    \qquad
    \eps_D=\epsinf-
    \frac{\omegap^2}{\omega(\omega+\ii\gamma)}.
    \label{eq:model-parameters-summary}
\end{equation}
The coefficient \(\beta\) is chosen according to the response being matched;
the common collisionless optical choice is
\(\beta^2=3\vF^2/5\), while static compressibility gives
\(\beta^2=\vF^2/3\).

\subsection*{What is and is not included}

Vanilla hard-wall HDM includes finite compressibility, nonlocal longitudinal
response, and a no-flux electron boundary condition.  It does not include:
\begin{itemize}
    \item equilibrium-density spill-out;
    \item tunnelling through a dielectric gap;
    \item explicit exchange--correlation potentials;
    \item microscopic surface scattering or Landau damping; or
    \item a spatially varying equilibrium density.
\end{itemize}
Those effects require extensions such as self-consistent quantum hydrodynamics,
the generalized nonlocal optical response model, quantum-corrected models, or
surface-response/Feibelman-parameter descriptions.

\begin{thebibliography}{9}

\bibitem{Ritchie1957}
R.~H. Ritchie,
\newblock Plasma losses by fast electrons in thin films,
\newblock \emph{Physical Review} \textbf{106}, 874 (1957).

\bibitem{Boardman1982}
A.~D. Boardman, ed.,
\newblock \emph{Electromagnetic Surface Modes},
\newblock Wiley, Chichester (1982).

\bibitem{Halevi1995}
P. Halevi,
\newblock Hydrodynamic model for the study of optical nonlocalities in metals,
\newblock \emph{Physical Review B} \textbf{51}, 7497 (1995).

\bibitem{Mermin1970}
N.~D. Mermin,
\newblock Lindhard dielectric function in the relaxation-time approximation,
\newblock \emph{Physical Review B} \textbf{1}, 2362 (1970).

\bibitem{Mortensen2014}
N.~A. Mortensen, S.~Raza, M. Wubs, T. S\o rensen, and S. Bozhevolnyi,
\newblock A generalized non-local optical response theory for plasmonic nanostructures,
\newblock \emph{Nature Communications} \textbf{5}, 3809 (2014).

\end{thebibliography}

\end{document}
